\scanpage{177}\mspage{32}%
% C original page filename: H134-05_134-2_scan177.tex
% Main writing: letter, his hand. Author's number 32 is at the upper edge.
où ``pris à la lettre'' $p''$ est maintenant une inclusion, mais ``moralement'' c'est \inserted{tjrs} une fibration à fibres très \unsure{polies|jolies} [de dim finie notamment, voire des variétés non singulières - c'est donc beaucoup mieux que ce que nous donne le yoga de Cartan-Serre ``toute application continue est équivalente à une fibration''\ldots] C'est ce dernier diagramme qui est susceptible d'une construction directe, purement algébrique, dans le cadre des schémas, lorsque l'on a développé la construction du voisinage tubulaire étale.

Le point où je veux en venir, c'est que la considération \inserted{de la cat} des faisceaux (ou \inserted{complexes d'iceux, ou} $n$-champs\unsure{,} \unsure{..|\ldots}) $\unsure{\mathcal P|P}$-constructibles sur un $X$, où $\mathcal P$ est une stratification ``équisingulière'' donnée, revient dans le cas particulier actuel à la connaissance d'un diagramme de \uline{types d'homotopie}%
% Isolated short nontext mark near (1633,1539), right of the underlined line; origin uncertain.
\raisebox{0pt}[0pt][0pt]{\parbox{0pt}{\figure{\begin{tikzpicture}[overlay,x=1pt,y=1pt]\draw[line width=0.4pt] (22,-2) -- (22.2,-0.7);\end{tikzpicture}}}} ordinaires \inserted{(ou de pro-types, si on revient à la top. étale)}%
% The parenthesis is interlinear above ordinary-baseline ordinaires. No guide identifies ordinaires as inserted.
% Separate short nontext stroke near (1334,1565), after and above the closing parenthesis; origin uncertain.
\raisebox{0pt}[0pt][0pt]{\parbox{0pt}{\figure{\begin{tikzpicture}[overlay,x=1pt,y=1pt]\draw[line width=0.3pt] (6,4) -- (6.2,5.5);\end{tikzpicture}}}}
\diagram{\begin{tikzcd}[ampersand replacement=\&,column sep=large,row sep=large]
\unsure{\Delta|\triangle} \arrow[r,"\unsure{J|j}"] \arrow[d,"p"'] \& X_1\\
X_0 \&
\end{tikzcd}}
en %
% Isolated short crooked nontext mark near (315,1848), above prenant; origin uncertain.
\raisebox{0pt}[0pt][0pt]{\parbox{0pt}{\figure{\begin{tikzpicture}[overlay,x=1pt,y=1pt]\draw[line width=0.45pt] (2,13) .. controls (2.7,13.5) and (1.8,14) .. (2,14.5);\end{tikzpicture}}}}prenant des systèmes de coeff locaux (ou de $n$-champs loc. constants) sur ces sommets \inserted{$X_0,X_1$ \deleted{$\unsure{\Delta|X}$}} reliés par des hom de compatibilité du type $p^*(F_0)\to J^*\unsure{(}F_1)$. Ce serait un exercice amusant (que je n'\unsure{ai} pas encore fait) de
% archive sheet number: 177
